\chapter{Professional Data Engineer Exam Mapping}
\index{exam!Professional Data Engineer}\index{certification}\index{study strategy}%
This appendix maps the Google \emph{Professional Data Engineer} certification onto the book's
chapters, so the final study weeks can be run as a targeted loop rather than a reread. It expands
the exam-preparation Friday of Chapter~24 \cite{googlecloudpdeexam}.

\section{Exam Domains Mapped to Chapters}
The exam assesses design and operational judgment, not syntax recall. The official guide groups
objectives into domains; the table below assigns each domain its chapters, the specific objectives
the book covers, and the trap the domain sets for unprepared candidates.

\begin{center}
\footnotesize
\begin{tabular}{@{}p{3.6cm}p{2.4cm}p{6.2cm}@{}}
\toprule
\textbf{Exam domain} & \textbf{Chapters} & \textbf{Objectives covered in this book} \\
\midrule
Designing data processing systems & 1--4, 21--23 & Storage class and location selection; relational vs.\ Spanner vs.\ Bigtable vs.\ BigQuery selection; migration and HA design; security and compliance baked into design \\
Ingesting and processing data & 9--15 & Batch (Dataproc, Data Fusion, Dataflow batch), CDC (Datastream), orchestration (Composer), messaging (Pub/Sub), streaming (Dataflow windowing, state, exactly-once) \\
Storing the data & 1--8 & Object storage lifecycle; warehouse partitioning, clustering, materialized views; external/federated storage; lakehouse (BigLake, Iceberg) \\
Preparing and using data for analysis & 5--8, 16--17 & SQL depth, modeling (grain, SCD2), BigLake governance, sharing (Analytics Hub), BI and semantic layers (Looker), in-warehouse ML (BQML) \\
Maintaining and automating data workloads & 12, 18--20, 24 & Idempotent orchestration, quality gates (assertions, AutoDQ), MLOps (pipelines, registry, monitoring), DataOps (Terraform, Dataform, CI/CD) \\
Security, compliance, governance & 21--23 & IAM least privilege, VPC Service Controls, CMEK/CSEK, Dataplex catalog/lineage/quality, DLP discovery and de-identification, erasure design \\
\bottomrule
\end{tabular}
\end{center}

\noindent\textbf{Typical trap per domain} (from Chapter~24): self-managed answers where a managed
service exists; batch answers to streaming requirements; ignoring lifecycle and partition cost
levers; skipping governance and semantics; solutions with no monitoring or CI/CD; IAM-only answers
that ignore perimeters.

\section{The Case Studies}
\index{case studies}%
The exam's case studies present fixed business scenarios. \textbf{EHR Healthcare} is a
compliance-driven migration: the Chapter~11 CDC patterns, Chapter~21 perimeter, and Chapter~23 DLP
designs answer most of its questions. \textbf{Helium Sports} is real-time analytics at scale:
Chapters~13--16 streaming patterns. Read both case-study documents before the exam and annotate
each requirement with the GCP service and this book's chapter that answers it; most case-study
questions reduce to a selection decision already made hands-on.

\section{Study Strategy}
\begin{enumerate}
    \item \textbf{Baseline mock (week -3).} One full-length practice exam, closed book, untimed
    but tracked. Score \emph{by domain}, not in aggregate---a 60\% in one domain matters more
    than an 82\% average.
    \item \textbf{Targeted rebuild loop (weeks -3 to -1).} For each weak domain, rebuild that
    week's Thursday lab from the solution sketches in the lab-solutions appendix, re-read the
    Friday use case, and re-answer the chapter's self-checks. Hands-on rebuilds beat rereading.
    \item \textbf{Timed mock (week -1).} Second full mock under time pressure. The
    schedule-the-real-exam threshold: above 85\% overall with no domain below 70\%.
    \item \textbf{Exam day heuristics.} Identify the stem's stated constraint (cost, latency,
    compliance, operational simplicity) first; eliminate the two answers that ignore it; prefer
    managed over self-managed, least-privilege over convenient, and the answer that includes
    monitoring over the one that merely works.
\end{enumerate}

\section{Self-Assessment Checklist}
Tick a box only when you could perform the task hands-on, unprompted, in the console or CLI.

\subsection*{Design (Chapters 1--4, 21)}
\begin{itemize}
    \item[$\square$] Choose storage class, location, and lifecycle for a stated access pattern and budget.
    \item[$\square$] Choose among Cloud SQL, AlloyDB, Spanner, Bigtable, and BigQuery for a stated workload, naming the deciding requirement.
    \item[$\square$] Design migration (DMS, Datastream, Transfer Service) with rollback and validation.
    \item[$\square$] Sketch the security envelope (IAM, perimeter, CMEK) before the data design.
\end{itemize}

\subsection*{Ingestion and Processing (Chapters 9--15)}
\begin{itemize}
    \item[$\square$] Build and tear down an ephemeral Dataproc job; explain preemptible trade-offs.
    \item[$\square$] Author a Data Fusion pipeline with Wrangler directives and macros.
    \item[$\square$] Stand up a Datastream CDC pipeline with backfill and a contract test.
    \item[$\square$] Write an idempotent, sensor-triggered Composer DAG with an \texttt{all\_done} notifier.
    \item[$\square$] Configure Pub/Sub schemas, dead-lettering, and replay.
    \item[$\square$] Implement windows, watermarks, triggers, and bounded state in Beam on Dataflow.
\end{itemize}

\subsection*{Storage and Analysis (Chapters 5--8, 16)}
\begin{itemize}
    \item[$\square$] Partition and cluster a table and prove the bytes-billed reduction.
    \item[$\square$] Decide on-demand versus Editions from measured slot usage.
    \item[$\square$] Model a star schema with declared grain and an SCD2 dimension.
    \item[$\square$] Govern lake files with BigLake policy tags and row access policies; share via Analytics Hub.
    \item[$\square$] Define one metric in a semantic layer and consume it in a dashboard.
\end{itemize}

\subsection*{ML and Operations (Chapters 17--20, 24)}
\begin{itemize}
    \item[$\square$] Train, evaluate, and serve a BQML model with leakage-free features.
    \item[$\square$] Deploy a Vertex model behind an endpoint with registry lineage and a rollback path.
    \item[$\square$] Gate automated retraining on drift evidence with quality preconditions.
    \item[$\square$] Wrap external AI API calls in batching, backoff, and dead-lettering, with cost forecast.
    \item[$\square$] Provision with Terraform from locked remote state; gate deploys through CI with Workload Identity Federation.
\end{itemize}

\subsection*{Security and Governance (Chapters 21--23)}
\begin{itemize}
    \item[$\square$] Enclose a project in a VPC SC perimeter and prove it with a negative test.
    \item[$\square$] Enforce column-level security with policy tags and row access policies, both negative-tested.
    \item[$\square$] Run DLP discovery and de-identification, including FPE tokenization with a KMS-wrapped key.
    \item[$\square$] Explain erasure design across time travel, backups, and replicas.
    \item[$\square$] Trace lineage from a dashboard tile to its source system in Dataplex.
\end{itemize}

\begin{tipbox}
If a box cannot be ticked honestly, that line \emph{is} the study plan: find the chapter, rebuild
the lab, and re-attempt the mock for that domain. The exam is passable on judgment built from
these labs; it is not passable on vocabulary alone.
\end{tipbox}
